\documentclass{article}
\usepackage{graphicx}

\title{Research report thesis}
\author{Mano van Holten}
\date{November 2025}

\begin{document}

\maketitle

\section{Instructions}

\begin{enumerate}
    \item Abstract (not mandatory)
    \item Write like a research paper
    \item No appendices
    \item Min 1, Max 4 figures and tables 
    \item 2500 words, figures and tables do not count
    \item introduction (sota), methods, discussion
\end{enumerate}

\section{Introduction}
\begin{itemize}
    \item copy from proposal and add / elaborate on:
    \item What is a DAG (in figure title).
    \item Think about this: we first regress on the covariates and then use residuals. The resulting regression uses only one (or two) degrees of freedom. In reality, there is uncertainty about the estimates of the covariate effects, which should be taken into account! See the manuscript + code for the bootstrapping approach!
    \item How does this interact with the chosen machine learning technique?
    \item Specify which machine learning techniques are chosen, why, and provide a place to learn more about them.
    \item Explain oversmoothing and overfitting and the bias-variance trade-off. Also get a more complete picture about plug-in bias. IS plug-in bias actually this underestimation? Also, what is meant with 'naive' applications?
\end{itemize}

\section{Methods}
\begin{itemize}
    \item Describe data simulation based on matrix algebra.
    \item Explain the chosen sample sizes based on research in the social sciences: what are plausible choices? How does this interact with the choice of machine learning method: think about the capacity of a model vs the ammount of data we have. There are guidelines of ML methods with sample sizes. Also think about the sparsity of the dataframes in high-dimensional space. We do not really have this problem since we have only X and Y: hence we might be able to get away with a slightly more complicated algorithm. 
    \item Explain the chosen types of non-linearity. Some realistic and some stress tests. Also explain what is meant with the extent of non-linearity: can we think of a measure that quantifies how non-linear the relationship is: i.e. the error left-over when fitting a linear regression - the added variance? This obviously is dependent on the number of non-linear terms and their chosen parameter values. 
    \item Explain the chosen parameter values. Like the innovations, cross-lagged and auto regressive parameters and also the left-over correlation of the innovations (0.3). 
    \item Specify precisely which cross-lagged models will be chosen and provide an overview of their effect sizes. Also provide expectations on how well these models deal with the non-linear behavior.  
    \item Specify which machine learning techniques are chosen, why, and provide a place to learn more about them. I was thinking about Tab-PFN since its SOTA on small tabular datasets, and its pre-training will avoid overfitting. In a sense this balances complexity and fit while not so reliant on the precise data, which might be very desirable for practitioners since they do not have to deal with any hyperparameter nonsense. TabPFN is said to have similar performance to a very well tuned XGboost algorithm, but do applied researchers really want to deal with that? Maybe also choose a very simple machine learning method that encapsulates the non-linearity generated by social processes well, such as the lasso, ridge or elnet. The nice thing about these is that we can work with "up to" cubic terms and interactions, and the algorithms requires very little tuning except for the lambda value, which we could specify (fix) at a specified 'general' value for social science research which would avoid meta-overfitting and be great for benchmarking. Of course, we would need to be able to show that this is 'enough' complexity, such that xgboost or RF is not needed, or at least difficult to train on limited data.
\end{itemize}



\section{Results}


\section{Discussion}


\end{document}
